\section{Tétlenség Ára}

Egyesek hisznek még a képviseleti demokráciában, és úgy érzik, hogy elegendő a
következő választáson a megfelelő képviselőnek mandátumot adni. Mások feladták,
és nem érzik szükségességét, hogy bármilyen formában részt vegyenek a politikában;
szerintük a dolgok előbb-utóbb rendbe jönnek, elég, ha valaki csupán magára koncentrál.

Ebben a fejezetben azt kívánjuk bemutatni, hogy mindkét csoport téved, és ez tényadatokkal
kimutatható.  Azok a krízisek, amelyek a közember mindennapjait aktívan rontják,
nem újak. Legyen szó lakhatási, népességi, gazdasági vagy ökológiai válságról,
egyértelmű lineáris romlás mutatható ki, évtizedes történettel.  Valószínűtlen, hogy
ezen válságok javulni fognak rendszerszintű változás nélkül. Akármilyen határozottan
húzzuk be az X-et a következő választáskor, vagy akármilyen mélyre is dugjuk a fejünket
a homokba, a tény megkerülhetetlen: vagy változunk, vagy egy sokkal üresebb, szürkébb,
és mindent egybevéve rosszabb élet vár ránk az eljövendő évtizedekben.

Először végigmegyünk a legfontosabb válságokon röviden, és megmutatjuk a lineáris
romlási trendet, majd felvázolunk egy lehetséges jövőt, amely vár ránk, ha nem lépünk.


\subsection{Ökológiai Válság}

A globális felmelegedés hatásait ma már mindenki a bőrén érzi. Az 1960-as évek óta
egyértelmű növekedés látható a bolygó átlagos felszíni hőmérsékletében. Ennek
eredménye a kiszámíthatatlan időjárás, a minden évben hőmérsékleti rekordot döntő
nyarak, az Alföld elsivatagosodása, az aszályok és megannyi környezeti katasztrófa. Ezalatt
a 60 év alatt kormányok jöttek-mentek minden demokráciában, ennek ellenére egy
ország sincs közelebb az olajfüggőség megszüntetéséhez. Annak ellenére, hogy az
alternatívák elérhetők és olcsóbbak, mint a fosszilis energiahordozók \cite{renewables_1,
renewables_2}. Az, hogy nem léptünk ez ügyben, egy döntés volt: egy döntés, amelyet nem
a köznép, hanem néhány elit hozott, akik profitálnak a tétlenségünkből.

Az ijesztő dolog nem csupán az, hogy a hőmérséklet egyre gyorsabban emelkedik, hanem az is,
hogy nem tudjuk pontosan, hol van az a pont, amely után már nincs visszaút. Ha nem is érjük el
ezt a pontot hamarosan --- ami egyre valószínűtlenebb ---, akkor is a következő 5-10 évben
a nyarak olyan melegek lehetnek, hogy napközben egyszerűen nem lehet kint tartózkodni. Ez azt
jelenti, hogy egy, maximum két országgyűlési választásunk van arra, hogy felkészüljünk
erre az új valóságra.

Meddig emelkedhet a hőmérséklet, amíg a döntséhozók komolyan veszik a problémát? Tényleg
elég, ha egyszer elmegyünk szavazni, vagy egyáltalán nem szavazunk, ha meg akarjuk fordítani
ezt a trendet?

\insertPic{world_temp}{0.5}{Az ábra a globális átlag felszíni hőmérsékletet mutatja 1850
	és 2024 között \cite{world_temp_1}. A vörös vonal jelzi a trendet, amely egyértelműen
	egyre gyorsuló melegedést mutat.  Például 2024 nyarán az átlaghőmérséklet már elérte
	a \(1.6\) \textdegree C-ot. Ez a hőmérséklet a teherbírásuk határáig feszít több az
	ember által alkotott rendszert, mint például villamoshálózatot. Azonban valószínűtlen,
	hogy a melegedés megáll ezen a ponton. Szakértők \(2 - 2.5\) \textdegree C-ot jósolnak
	2050-re \cite{world_temp_2}, ami katasztrofális lenne.}


\subsection{Megélhetési Válság}

A fizetések stagnálása és az árak emelkedése látható mindenki által. Ez a probléma is
évtizedek óta fennáll, de az utóbbi években a helyzet itt is gyorsuló romlást mutat. Ugyan
az itt bemutatott adatok az Egyesült Államokat elemzik, minden országban hasonló a helyzet.

A számadatokon egyértelműen látszik, mi a probléma. Mindenki többet dolgozik kevesebb
bérért, sokaknak több állásuk is van. Azonban a kormányzati szakértők és a nagyobb
médiacégek csak néhány gazdasági mutatóval foglalkoznak, mint a GDP, átlagfizetések
vagy a tőzsdeindexek --- azonban ezek félrevezető mutatók.

Például, a GDP (bruttó hazai termék) egy gazdasági mutató, amely egy országban egy adott
időszak alatt előállított végső termékek és szolgáltatások összértékét méri. A
modern GDP-koncepciót az 1930-as években dolgozták ki, és a második világháború idején
fontos eszközzé vált a gazdasági tervezésben. Azonban megalkotója maga is figyelmeztetett
arra, hogy a GDP nem alkalmas egy nemzet jólétének mérésére, mivel nem veszi figyelembe
a jövedelmi egyenlőtlenségeket, a környezeti károkat, a nem fizetett munkát vagy az
életminőség más fontos tényezőit. Egy akkumulátorgyár, amely megmérgezi a környéken
élőket, ugyanolyan pozitívan fog a kalkulációba kerülni, mint egy cég, amelynek pozitív
hatása van a társadalomra.

Ugyanígy, az átlagbér önmagában nem tekinthető megbízható indikátornak a lakosság
pénzügyi jólétének értékelésére, mivel érzékeny a jövedelemeloszlás aszimmetriájára
és a szélsőértékekre. Erősen jobbra ferde jövedelemeloszlás esetén néhány magas
jövedelmű egyén jelentősen megemelheti az átlagot, miközben a mediánjövedelem és a
lakosság nagy részének tényleges anyagi helyzete változatlan marad. Emiatt az átlagbér
torzíthatja a gazdasági jólét percepcióját, és nem tükrözi megfelelően a jövedelmi
egyenlőtlenségeket vagy a relatív szegénység mértékét. A pénzügyi jólét pontosabb
méréséhez ezért célszerű mediánjövedelmet, jövedelmi kvantiliseket és egyenlőtlenségi
mutatókat --- például Gini-indexet --- is alkalmazni.

Végül pedig, a tőzsde teljesítménye nem tekinthető a gazdaság egészségének vagy
az átlagos háztartások jólétének megbízható indikátoraként, mivel a részvénypiac
elsősorban a vállalati profitvárakozásokat, a likviditási feltételeket és a befektetői
hangulatot tükrözi, nem pedig a reálgazdaság minden szegmensének állapotát. A
részvényindexek erősen koncentráltak néhány nagy kapitalizációjú vállalatra, így a
piaci emelkedések gyakran nem reprezentálják a kis- és középvállalatok teljesítményét
vagy a munkaerőpiaci helyzetet. Emellett a pénzügyi eszközök árfolyamai elszakadhatnak
a reáljövedelmektől, különösen monetáris lazítás és eszközár-infláció esetén,
ami tovább gyengíti a kapcsolatot a tőzsdei "bull market" és a háztartások tényleges
életszínvonala között. Következésképpen a gazdasági jólét értékeléséhez komplexebb
indikátorrendszer szükséges, amely magában foglalja a foglalkoztatottságot, a reálbéreket
és a jövedelmi egyenlőtlenséget is.

Jól látható tehát, hogy a GDP, az átlagbér, vagy a tőzsdei teljesítmény nem alkalmas
arra, hogy leírja egy gazdaság állapotát.  Ennek ellenére ezek a mutatók kapják a legtöbb
figyelmet a híradásokban. A gazdasági teljesítmény és az órabérek elválása (\cref{fig:wage_vs_prod})
minden nyugati országban látható, hol szélesebb, hol kevésbé széles a szakadék a két mutató között.
Ennek a jelenségnek az oka rendkívül összetett, amelyet nem tudunk megvitatni itt, de az
érdeklődők számára a következő referenciák érdekesek lehetnek: \cite{wage_2, wage_3}.

\insertPic{wage_vs_prod}{0.4}{Az ábra a gazdasági teljesítmény és az órabér erős
	elválását mutatja \cite{wage_1} az Egyesült Államokban.  Jól látható, hogy 1973 óta a
	divergencia nő. 1973-2013 között a gazdasági teljesítmény \(74.4\)\%-kal, míg az
	órabérek csupán \(9.2\)\%-kal nőttek. Ezt az átlag munkavállaló stagnáló fizetésként
	érzékeli.}


\subsection{Lakhatási Válság}

A lakhatási helyzet tarthatatlanságát mindenki jól érzi, aki 2000 után született. Hosszú
évtizedek alatt megannyi kormány, megannyi országban próbálta orvosolni a problémát,
de sikertelenül. Csakúgy mint a többi válság esetén, nem azért nem született megoldás
még, mert az ne létezne, hanem mert a magas ingatlanárakból profitálnak a tehetősebbek,
egy csoport, amely túlképviselt a döntéshozásban.

\insertPic{price_to_wage}{0.75}{Az ábra \cite{housing_2} az Egyesült Államok lakhatási
	megfizethetőségi indexét mutatja. Ez a szám egyenlő a medián házár osztva a medián
	háztartási jövedelemmel. Látható, hogy a helyzet hasonló a 2008-as ingatlanbuborékhoz.}

A megfizethetetlen lakhatás súlyos társadalmi problémákat okoz, mivel jelentősen korlátozza
a fiatalok életkezdési lehetőségeit, ezáltal fokozva az izolációt és a társas kapcsolatok
kialakításának nehézségét. A bérleti árak emelkedéséhez is hozzájárul, mivel azok
is kénytelenek bérlakásban maradni, akik egyébként ingatlanvásárlásra törekednének,
ez pedig további keresleti nyomást teremt a bérszektorban. Emellett növeli az egyének
kiszolgáltatottságát, például olyan helyzetekben, amikor a bántalmazó környezetből
való kilépés a lakhatás hiánya miatt ellehetetlenül. A jelenség erősíti a társadalmi
egyenlőtlenségeket is, mivel az ingatlantulajdonnal nem rendelkező csoportok jövedelmük
egyre nagyobb részét kénytelenek lakhatásra fordítani. Makroszinten a lakhatási válság
hozzájárul a társadalmi stabilitás gyengüléséhez, mivel növeli a munkaerőpiaci
függőséget és csökkenti az egyéni alkupozíciókat. Végső soron egy olyan strukturális
kockázati környezetet hoz létre, amelyben már kisebb anyagi vagy élethelyzeti sokkok is
a hajléktalansághoz vezethetnek.

\insertPic{avg_price_per_meter}{0.45}{Az ábra \cite{housing_1} a magyarországi használt
	ingatlanok négyzetméterárának átlagos értékét mutatja regiónként. Látható, hogy
	mindenhol egyértelmű növekedés figyelhető meg, Budapest esetében a leglátványosabb
	növekedés.}


\subsection{Politikai Válság}

A korábban említett válságok nem a véletlenek eredményei, és nem is csupán
ügyetlenségből erednek. Ezen krízisek olyan megfontolt intézkedések eredményei, amelyek
a politikában túlnyomórészt képviselt elitcsoportok érdekeit szolgálják. Ez magyarázza
azt, hogy miért nem segített a kormányváltás a megélhetési vagy lakhatási válságon az
USA-ban \cite{political_crisis_1, political_crisis_2}, Angliában \cite{political_crisis_3},
Németországban \cite{political_crisis_4, political_crisis_5}, sem Ausztráliában
\cite{political_crisis_6, political_crisis_7}.

Egy tanulmány \cite{political_crisis_8} empirikusan vizsgálta, hogy az amerikai politikában
kik gyakorolnak tényleges befolyást a közpolitikára: az átlagos állampolgárok, a gazdasági
elit vagy az érdekcsoportok. A szerzők 1779 politikai döntésen végzett statisztikai elemzés
alapján arra jutnak, hogy a gazdasági elit és az üzleti érdekcsoportok jelentős, önálló
hatással bírnak a politikai döntésekre, míg az átlagos állampolgárok preferenciái alig
vagy egyáltalán nem befolyásolják a kimeneteket. Ez a jelenség minden országban jelen
van, hol dominánsabban, hol kevésbé. Azonban a tény megkerülhetetlen: az érdekcsoportok
befolyása nagymértékben meghatározza a politikai irányt mindenhol.

Azonban a krízisek és ezek interszekciói már nem hagyhatók figyelmen kívül, és kifutunk az
időből. Nincs lehetőségünk tucatnyi választást kivárni egy olyan kormányzatért, amely
képes és hajlandó kezelni az átlagpolgárt érintő problémákat. Egy olyan rendszerre van
szükség, amely garantálni tudja a népakarat és a törvényhozás egybefonódását. Bármi
más csak tovább fogja rontani a helyzetet.


\subsection{Népességi Válság}

A korábban említett válságok gyökere, hogy a törvényhozás egy szűk, homogén elit
kezében van. Ez egyértelműen oda vezet, hogy ezek a csoportok a saját érdekeik mentén
hoznak döntéseket, ami egy olyan társadalmat eredményez, amely nem tudja reprodukálni
magát. Ahhoz, hogy a népesség stabil legyen, \(2.1\) gyermekre van szükség családonként.
Ezt Magyarország 1975 óta nem éri el. Azóta megannyi kormány próbálkozott, de nem tudta
megfordítani a trendet.

A népesség évről évre fogy, azonban ahelyett, hogy implementálnánk a fiatalok követeléseit
--- minden országban ---, milliomos idős férfiak azt próbálják kitalálni, hogyan vehetnék
rá a fiatalokat még több gyermek vállalására.

Különböző országok kormányai mindent megtettek, hogy biztosítsák a befektetők profitját
azáltal, hogy magasan tartották az ingatlanárakat, és segédprogramokat indítottak el,
hogy a családalapításra vágyók ki tudják fizetni az egyre növekvő ingatlanárakat. Ez
rövid távon segített, azonban a következő generációnak még nagyobb összegeket kell
kifizetnie. Úgy tűnik, a kormányok kifogytak a trükkökből, és most nehéz választásra
kényszerülnek: vagy lehozzák az ingatlanárakat, segítve a fiatalokat és ártva a
befektetőknek, vagy hagyják az árakat tovább emelkedni, megvédve a befektetőket, és
végignézik, ahogy lassan eltűnik az országuk.

\insertPic{birthrate}{0.5}{Az ábra \cite{birthrate_1} az egy nőre jutó átlagos gyermekek
	számát mutatja a World Bank adatai alapján. A népesség stabilitásához \(2.1\) gyermekre van
	szükség. Ezt az értéket az ország 1975 óta nem éri el, és azóta folyamatos csökkenés
	figyelhető meg. 2011-ben egy gyenge növekedés indult meg, azonban 2020-ra ez éles csökkenésre
	váltott. A legfrissebb adatok szerint az ország elérte az \(1.4\) átlagot, és várhatóan
	\(1.0\) körüli stabilizálódás várható. Ez azt jelentené, hogy a 8 milliós lakosság
	körülbelül 2 millióra fog zuhanni a következő 50 évben. Ez a becslés egy rendkívül
	egyszerűsített modell alapján készült, azonban függetlenül attól, hogy a tényleges
	szám \(1.5\) vagy \(3\) millió lesz, ez egy olyan jövő, amelyet érdemes lenne elkerülnünk.}


\subsection{Fényes Jövő?}

Ebben a fejezetben csak egy rövid előretekintést teszünk, amivel remélhetőleg az olvasók egy
része meggyőzhető, hogy a korábban említett krízisek nem oldhatóak meg tétlenséggel, vagy
ha a dolgokat a régi kerékvágásban hagyjuk folytatódni. Ha nem sikerül érdemi megoldásokat
találni ezekre a problémákra, akkor egy sokkal komorabb és üresebb jövő elé nézünk.

Természetesen nincs modell, amely pontos előrejelzést tudna adni a jövővel kapcsolatban,
de nem is ez a célunk. Sokkal inkább egy lehetséges jövőképet szeretnénk felvázolni
azon információk alapján, amelyek a rendelkezésünkre állnak, ezzel buzdítva az olvasót
cselekvésre.

\begin{displayquote} 2050-re a bolygó átlagos felszíni hőmérséklete elérte a \(2.1\)
	\textdegree C-ot. A nyarak elviselhetetlenül forróak. Mindenki légkondicionáló alatt
	gubbaszt otthon, és csak akkor hagyja el a házat, ha nincs más választása. Ez rendkívül
	megterheli az elektromos hálózatot, amelyet nem ilyen hőmérsékletekre terveztek. Gyakoriak
	az áramkimaradások, főleg a legmelegebb napokon.  A kiszámíthatatlan időjárás miatt a
	növénytermesztés nehézkes, sok a veszteség. Ez meg is látszik az élelmiszerárakban; az
	állati eredetű élelmiszerek ára magas, mivel a takarmány és az állattartás költségei
	is többszörösére nőttek az utóbbi évtizedekben. Ennek köszönhetően sokan nem vagy csak
	ritkán engedhetik meg maguknak, hogy húst tegyenek az asztalra. Minden nyáron egyre több
	a rosszullét napközben, ami miatt a kormány fontolgatja az éjszakai műszakok bevezetését.

	A fizetések sem alakultak jobban. A legtöbb ember kettő-három munkahelyen dolgozik, hogy
	lépést tudjon tartani a mindennapi szükségletek emelkedő áraival. Az adók is emelkednek,
	hiszen a költségvetési hiány egyre nagyobb, a közszolgáltatások fenntartása pedig egyre
	drágább. A lakbérek is jelentősen emelkedtek, köszönhetően a rekordmagas ingatlanpiacnak,
	amellyel a több ingatlannal rendelkezők jól jártak. A közbeszédben elterjed egy mondás,
	miszerint: "Egy állás a kormánynak, egy a főbérlőnek, és egy nekem", utalva az adók,
	lakbérek és egyéb költségek növekedésére.

	A saját otthon elérhetetlen vágyálom lett a legtöbbek számára. Azok, akik sikeresen
	hozzájutottak banki kölcsönhöz, 40-50 évig fogják nyögni a törlesztőrészletet. Sokan
	adósságot hagynak a gyermekeikre, de ők a szerencsések. Sokan olyan helyen vettek ingatlant,
	amely mára már elsivatagosodott, hajléktalanná téve őket. A lakásbiztosítás is sokat
	emelkedett, ami szintén sokakat elzárt az ingatlantulajdontól.

	Nagy az elégedetlenség a világban, ezért minden évben egy új, radikálisabb párt ígér
	változást, azonban a helyzet sosem változik. A közbiztonság jelentősen romlott, ami arra
	késztette a kormányt, hogy egy kiterjedt megfigyelőrendszert építsen ki. Üzenetek, képek,
	online tevékenységek mind monitorozva vannak.

	A világ kiszámíthatatlan helyzete és a gazdaság romlása miatt egyre kevesebben érzik
	felelős döntésnek a gyermekvállalást. Ez a statisztikákon is meglátszik: minden évben új
	negatív rekordot dönt a született gyermekek száma. Ez sok iskola és kórház bezárását
	jelentette, és még több terhet helyezett a nyugdíjrendszerre, amelybe egyre kevesebben tudnak
	befizetni. Egyes szakemberek szerint az egyetlen járható út, ha a nyugdíjkorhatárt tovább
	emeljük, azonban ez várhatóan még több elégedetlenséget fog szülni a társadalomban.
\end{displayquote}